\section{I2C Protocol Fundamentals}
\label{sec:i2c-fundamentals}

I2C is two wires and a clock, and everything else follows from one electrical decision: the
lines are \structure{open drain}. A device never drives a line high; it releases it and a
pull-up resistor does the work. Two consequences drive every design decision that follows.

\begin{keytakeaways}
    \item A device can only pull a line \emph{low}. ``High'' is the absence of anybody
          pulling.
    \item Because releasing is not a conflict, two devices can never damage each other by
          talking at once — the bus resolves it as a wired-AND.
    \item But because a ``1'' comes from a resistor, the rise time is set by
          $R \times C$: the bus capacitance is what limits the clock rate and the wire length.
\end{keytakeaways}

\subsection{The Transaction, in Voltage}

A transfer is a sequence of edges, and each of them means something specific:

\begin{enumerate}
    \item \textbf{START}: SDA falls while SCL is high. Nothing else on the bus produces that
          pattern, which is why it can be used as a marker without a separate wire.
    \item \textbf{Address}: seven bits, most significant first, then a direction bit
          ($0$ = write, $1$ = read).
    \item \textbf{ACK}: the addressed device pulls SDA low for the ninth clock. If nobody
          pulls, the byte is unacknowledged — the single most useful diagnostic on the bus.
    \item \textbf{Data}: bytes, each followed by an ACK, in the direction agreed.
    \item \textbf{STOP}: SDA rises while SCL is high.
\end{enumerate}

The detail that catches everybody the first time: \textbf{SDA may only change while SCL is
low}. Both the START and the STOP are defined as the violation of that rule, which is
exactly why they are unambiguous.

\subsection{The Electrical Reality on This Board}

The EFR32MG24 does not have a dedicated I2C pin pair in the sense a beginner expects. The
peripheral observes ordinary GPIO pins, so the same pin can be used as a GPIO, as an I2C
line, or as something else — and the routing is a separate configuration step.

That gives two ways to implement a bus on this part:

\begin{description}
    \item[The on-chip controller (I2C0/I2C1).] The hardware generates the edges, handles
          addressing, ACK and clock stretching, and interrupts the core when done.
    \item[A bit-banged master.] Software drives the pins. Slower, more code, and completely
          under your control — which matters enormously when the first option misbehaves.
\end{description}

This chapter uses both, and the reason is not stylistic. On this part the on-chip
controller could not complete a read, and the proof of that is the subject of the next
section.

\subsection{Open Drain in the Registers}

An open-drain pin is not configured ``as output and set to 1''. It is configured in a mode
where the output driver can only sink current, and the value written has the opposite
meaning to the one intuition suggests:

\begin{minted}{rust}
// The line is released (reads high through the pull-up) when the output latch is 1,
// and pulled low when it is 0. This is not a convention we chose: it is what WIREDAND
// means, and it is why a bit-banged master writes `1` to let go.
porta_dout().write(|w| unsafe { w.dout(1 << pin) }); // release
porta_dout().write(|w| unsafe { w.dout(0) });        // drive low
\end{minted}

Getting this backwards produces the most confusing failure in the whole chapter: a bus that
appears dead, while the code looks correct.

\begin{warningbox}
    \textbf{The mistake that cost the most time in this project.} The GPIO was configured for
    the bus, the address was right, the device was powered, and every scan returned nothing —
    because the pin was in a mode where writing \codeinline{1} \emph{drove the line high}
    instead of releasing it. Two open-drain lines both driven high is a short through the
    drivers, and the resulting waveform is not an I2C bus at all.
\end{warningbox}
